一维均匀介质体积热容 \(c_v>0\)、导热系数 \(\kappa>0\)、热流弛豫时间 \(\tau>0\)，温度是平衡态附近的小偏差 \(T(x,t)\)。采用能量连续性 \(c_vT_t+q_x=0\) 与Maxwell–Cattaneo关系 \(\tau q_t+q=-\kappa T_x\)。取周期边界或无限域中快速衰减的扰动，令 \(D=\kappa/c_v\)。
(1) 消去热流导出双曲型热传导方程，解释与瞬时Fourier扩散的差异和引入的初始数据。
(2) 令 \(T\propto e^{i(kx-\omega t)}\)，求两个复频分支和长波、短波极限，并计算热波传播速度和衰减率。
(3) 构造温度—热流联合正定二次泛函，证明精确局部能流恒等式与总量单调耗散；指出它不是原始热力学内能本身。
(4) 在固定 \(D\) 下取 \(\tau\to0^+\)，分别得到扩散慢模、快速弛豫模，解释为何短波 \(k\to\infty\) 与 \(\tau\to0\) 的极限顺序不可交换。
